\documentclass{amsart}
% For dealing with references we use the comment environment
\usepackage{verbatim}
\newenvironment{reference}{\comment}{\endcomment}
%\newenvironment{reference}{}{}
\newenvironment{slogan}{\comment}{\endcomment}
\newenvironment{history}{\comment}{\endcomment}

% For commutative diagrams we use Xy-pic
\usepackage[all]{xy}

% We use 2cell for 2-commutative diagrams.
\xyoption{2cell}
\UseAllTwocells

% We use multicol for the list of chapters between chapters
\usepackage{multicol}

% This is generally recommended for better output
\usepackage{lmodern}
\usepackage[T1]{fontenc}

% For cross-file-references

% Package for hypertext links:
\usepackage{hyperref}

% For any local file, say "hello.tex" you want to link to please

% Theorem environments.
%
\theoremstyle{plain}
\newtheorem{theorem}[subsection]{Theorem}
\newtheorem{proposition}[subsection]{Proposition}
\newtheorem{lemma}[subsection]{Lemma}

\theoremstyle{definition}
\newtheorem{definition}[subsection]{Definition}
\newtheorem{example}[subsection]{Example}
\newtheorem{exercise}[subsection]{Exercise}
\newtheorem{situation}[subsection]{Situation}

\theoremstyle{remark}
\newtheorem{remark}[subsection]{Remark}
\newtheorem{remarks}[subsection]{Remarks}

\numberwithin{equation}{subsection}

% Macros
%
\def\lim{\mathop{\mathrm{lim}}\nolimits}
\def\colim{\mathop{\mathrm{colim}}\nolimits}
\def\Spec{\mathop{\mathrm{Spec}}}
\def\Hom{\mathop{\mathrm{Hom}}\nolimits}
\def\Ext{\mathop{\mathrm{Ext}}\nolimits}
\def\SheafHom{\mathop{\mathcal{H}\!\mathit{om}}\nolimits}
\def\SheafExt{\mathop{\mathcal{E}\!\mathit{xt}}\nolimits}
\def\Sch{\mathit{Sch}}
\def\Mor{\mathop{\mathrm{Mor}}\nolimits}
\def\Ob{\mathop{\mathrm{Ob}}\nolimits}
\def\Sh{\mathop{\mathit{Sh}}\nolimits}
\def\NL{\mathop{N\!L}\nolimits}
\def\CH{\mathop{\mathrm{CH}}\nolimits}
\def\proetale{{pro\text{-}\acute{e}tale}}
\def\etale{{\acute{e}tale}}
\def\QCoh{\mathit{QCoh}}
\def\Ker{\mathop{\mathrm{Ker}}}
\def\Im{\mathop{\mathrm{Im}}}
\def\Coker{\mathop{\mathrm{Coker}}}
\def\Coim{\mathop{\mathrm{Coim}}}

% Boxtimes
%
\DeclareMathSymbol{\boxtimes}{\mathbin}{AMSa}{"02}

%
% Macros for moduli stacks/spaces
%
\def\QCohstack{\mathcal{QC}\!\mathit{oh}}
\def\Cohstack{\mathcal{C}\!\mathit{oh}}
\def\Spacesstack{\mathcal{S}\!\mathit{paces}}
\def\Quotfunctor{\mathrm{Quot}}
\def\Hilbfunctor{\mathrm{Hilb}}
\def\Curvesstack{\mathcal{C}\!\mathit{urves}}
\def\Polarizedstack{\mathcal{P}\!\mathit{olarized}}
\def\Complexesstack{\mathcal{C}\!\mathit{omplexes}}
% \Pic is the operator that assigns to X its picard group, usage \Pic(X)
% \Picardstack_{X/B} denotes the Picard stack of X over B
% \Picardfunctor_{X/B} denotes the Picard functor of X over B
\def\Pic{\mathop{\mathrm{Pic}}\nolimits}
\def\Picardstack{\mathcal{P}\!\mathit{ic}}
\def\Picardfunctor{\mathrm{Pic}}
\def\Deformationcategory{\mathcal{D}\!\mathit{ef}}

\setlength{\emergencystretch}{3em}
\begin{document}
\title[Stacks Project, Tag 07CZ]{341 --- Independence of polynomial presentation for Koszul-regular kernels}
\maketitle
\noindent Copyright (C) 2005--2025 Johan de Jong. Stacks Project authors. Source: \href{https://stacks.math.columbia.edu/tag/07CZ}{Tag 07CZ}.

\noindent Editorial difficulty note (not author text). Difficulty H2: The author constructs a common polynomial presentation and compares conormal dimensions of its two kernels. Local bases, concatenation of regular sequences, and a change-of-generators argument transfer Koszul regularity to the other presentation rather than assuming presentation independence..

\noindent Editorial provenance note (not author text). History: excerpt from revision \texttt{a0444\allowbreak 6e57e\allowbreak c1fbc\allowbreak 252a8\allowbreak 71afc\allowbreak ec775\allowbreak 2fb28\allowbreak 07b14}, more-algebra.tex, lines 8259--8332. Reference presentation was adapted; original-author context and core support blocks were retained or added. The mathematical wording is unchanged. Licensed under GNU FDL 1.2 or later; no invariant sections or cover texts. See ../licenses/COPYING. Context blocks reproduce author definitions and supporting results; they are not additional counted problems.

\begin{lemma}
\label{lemma-koszul-independence-presentation}
Let $A \to B$ be a finite type ring map. If for some presentation
$\alpha : A[x_1, \ldots, x_n] \to B$ the kernel $I$ is a Koszul-regular ideal
then for any presentation $\beta : A[y_1, \ldots, y_m] \to B$ the kernel
$J$ is a Koszul-regular ideal.
\end{lemma}

\begin{proof}
Choose $f_j \in A[x_1, \ldots, x_n]$ with $\alpha(f_j) = \beta(y_j)$
and $g_i \in A[y_1, \ldots, y_m]$ with $\beta(g_i) = \alpha(x_i)$.
Then we get a commutative diagram
$$
\xymatrix{
A[x_1, \ldots, x_n, y_1, \ldots, y_m]
\ar[d]^{x_i \mapsto g_i} \ar[rr]_-{y_j \mapsto f_j} & &
A[x_1, \ldots, x_n] \ar[d] \\
A[y_1, \ldots, y_m] \ar[rr] & & B
}
$$
Note that the kernel $K$ of $A[x_i, y_j] \to B$ is equal to
$K = (I, y_j - f_j) = (J, x_i - g_i)$. In particular, as
$I$ is finitely generated by
Lemma \href{https://stacks.math.columbia.edu/tag/07CW}{lemma quasi regular ideal finite (Tag 07CW)}
we see that $J = K/(x_i - g_i)$ is finitely generated too.

\medskip\noindent
Pick a prime $\mathfrak q \subset B$. Since
$I/I^2 \oplus B^{\oplus m} = J/J^2 \oplus B^{\oplus n}$
(Algebra, Lemma \href{https://stacks.math.columbia.edu/tag/00S5}{lemma conormal module (Tag 00S5)})
we see that
$$
\dim J/J^2 \otimes_B \kappa(\mathfrak q) + n =
\dim I/I^2 \otimes_B \kappa(\mathfrak q) + m.
$$
Pick $p_1, \ldots, p_t \in I$ which map to a basis of
$I/I^2 \otimes \kappa(\mathfrak q) = I \otimes_{A[x_i]} \kappa(\mathfrak q)$.
Pick $q_1, \ldots, q_s \in J$ which map to a basis of
$J/J^2 \otimes \kappa(\mathfrak q) = J \otimes_{A[y_j]} \kappa(\mathfrak q)$.
So $s + n = t + m$. By Nakayama's lemma there exist $h \in A[x_i]$ and
$h' \in A[y_j]$ both mapping to a nonzero element of $\kappa(\mathfrak q)$
such that $I_h = (p_1, \ldots, p_t)$ in $A[x_i, 1/h]$ and
$J_{h'} = (q_1, \ldots, q_s)$ in $A[y_j, 1/h']$.
As $I$ is Koszul-regular we may also assume that $I_h$ is generated
by a Koszul regular sequence. This sequence must necessarily have length
$t = \dim I/I^2 \otimes_B \kappa(\mathfrak q)$, hence we see that
$p_1, \ldots, p_t$ is a Koszul-regular sequence by
Lemma \href{https://stacks.math.columbia.edu/tag/066A}{lemma independence of generators (Tag 066A)}.
As also $y_1 - f_1, \ldots, y_m - f_m$ is a regular sequence we
conclude
$$
y_1 - f_1, \ldots, y_m - f_m, p_1, \ldots, p_t
$$
is a Koszul-regular sequence in $A[x_i, y_j, 1/h]$
(see Lemma \href{https://stacks.math.columbia.edu/tag/0669}{lemma join koszul regular sequences (Tag 0669)}).
This sequence generates the ideal $K_h$. Hence the
ideal $K_{hh'}$ is generated by a Koszul-regular sequence
of length $m + t = n + s$. But it is also generated by the sequence
$$
x_1 - g_1, \ldots, x_n - g_n, q_1, \ldots, q_s
$$
of the same length which is thus a Koszul-regular sequence by
Lemma \href{https://stacks.math.columbia.edu/tag/066A}{lemma independence of generators (Tag 066A)}.
Finally, by
Lemma \href{https://stacks.math.columbia.edu/tag/068M}{lemma truncate koszul regular (Tag 068M)}
we conclude that the images of $q_1, \ldots, q_s$ in
$$
A[x_i, y_j, 1/hh']/(x_1 - g_1, \ldots, x_n - g_n)
\cong A[y_j, 1/h'']
$$
form a Koszul-regular sequence generating $J_{h''}$. Since $h''$
is the image of $hh'$ it doesn't map to zero in $\kappa(\mathfrak q)$
and we win.
\end{proof}
\end{document}
